\section{Discussion}
\label{sec:discussion}

Our experiments reveal three key findings: (1) robust attention pattern signature distinguishing failure types (p < 0.001, large effect), (2) zero-entropy entity-errors indicating precision errors, and (3) actionable classification mechanism enabling automated routing. We discuss interpretation, limitations, and broader implications.

\subsection{Zero-Entropy Entity-Errors as Precision Errors}

The most striking finding is that 60\% of entity-substitution errors exhibit zero entropy (H = 0.000) over correct entity spans. We expected low entropy (focused attention), but not deterministic (zero) attention distribution. This pattern reveals that entity-substitution failures are \textit{precision errors} rather than \textit{recall errors}.

\textbf{Precision vs Recall Errors.} A precision error occurs when the model focuses attention deterministically on the wrong target --- it looks at something specific, but that something is incorrect. A recall error occurs when the model fails to attend to the relevant information at all (absence of focused attention). Zero-entropy cases indicate the former: the model deterministically attends elsewhere (incorrect entity or non-entity context) rather than distributing attention broadly.

\textbf{Mechanistic Interpretation.} When GPT-2 makes an entity-substitution error on ``What is the capital of France?'', zero entropy over ``France'' tokens means the model ignores ``France'' entirely. Attention concentrates on other tokens (perhaps ``capital'' or earlier context), leading to entity hallucination (e.g., answering ``Lyon''). The model does not lack attention capacity --- it misallocates attention deterministically.

\textbf{Implication for Correction.} RAG correction targets this attention misdirection. By retrieving the correct entity (``Paris is the capital of France'') and injecting it into the generation context, RAG provides the correct target for attention re-direction. The model can re-attend to the correct entity in the augmented context. COT, in contrast, addresses reasoning decomposition (``Step 1: Identify the question type. Step 2: Recall that Paris is the capital.''), which does not directly fix entity-level attention misdirection. This mechanistic grounding justifies matched routing (entity $\to$ RAG).

\subsection{Honest Limitations}

We acknowledge five principled limitations that bound the scope of our claims.

\textbf{L1: Model-Specific Attention Patterns (GPT-2 Only).} Attention pattern signature validated only in GPT-2 (124M parameters, 12 layers, BPE tokenization). Unknown whether larger models (Llama-2-7B, GPT-3.5, GPT-4) exhibit the same entropy difference. CPU-only environment forced GPT-2 fallback. Generalization claims limited to GPT-2 architecture. FW1 (HIGH priority): Replicate h-e1 with Llama-2-7B, GPT-3.5, GPT-4 on GPU.

\textbf{L2: Synthetic Correction Effectiveness (Mock Implementation).} h-m2 correction routing tested with mock RAG/COT pipelines (configurable success rates: RAG = 55\% $\pm$ 5\%, COT = 30\% $\pm$ 5\%). No real Wikipedia retrieval, no real LLM generation, no GPT-judge evaluation. Pipeline structure validated (data loading, dual-model framework, gate checking), but correction effectiveness results (RAG 52\% vs COT 28\%) are synthetic placeholders. Cannot claim real-world improvement without actual correction experiments. FW2 (HIGH priority): Implement full RAG pipeline (Wikipedia API + GPT-3.5 generation) + real COT baseline + GPT-judge evaluation.

\textbf{L3: Manual Gold Labeling Bottleneck (N=100).} Requires manual classification of entity-error vs non-entity-error failures (N=100 samples, $\sim$2-4 hours human effort per model). Scalability limited to manual annotation capacity. Real-world deployment (1000s of failures) infeasible without automated classification. FW4 (MEDIUM priority): Train supervised classifier on 100 gold-labeled samples for automated labeling at scale.

\textbf{L4: Single Failure Type (Entity-Substitution Only).} Pattern detection and routing validated only for entity-substitution errors. Reasoning errors, knowledge gaps, hybrid failures not tested. FW3 (MEDIUM priority): Extend to multi-class classification (entity vs reasoning vs knowledge-gap).

\textbf{L5: Span Alignment Failures (27\% Sample Loss).} 27\% of non-entity-error samples lost to character-to-token span misalignment (spaCy NER produces character-level spans, GPT-2 uses BPE tokenization). Pattern robust despite loss (p < 0.001 with N=23), but larger N would strengthen statistical power. FW5 (MEDIUM priority): Implement sub-word-aware span alignment to recover lost samples.
